%======================================================================
\subsection{SIM-06 --- Built, 31 Oct: the flow in the 3D view (and the FemPostPipeline crash)}\label{sec:sim06-post-built}
%======================================================================

\paragraph{What you can do now.} \textbf{Show in 3D} on an airflow study converts the last time with foamToVTK and
puts the result in the 3D view inside the app: the \emph{surface pressure} on each part (a pipeline of its boundary
patch, coloured by $p$), and the \emph{mid cut plane} (XZ) coloured by the velocity magnitude $|U|$; the volume
pipeline itself is added to the study but hidden, because it would fill the view. The far-field patches (inlet,
outlet, sides) are left out. Pressure and velocity stay in SI; only the points are scaled to millimetres so the flow
sits on the model. \textbf{Streamlines} stay in ParaView: FreeCAD's FEM module has no streamline filter object in
this build, so there is nothing to call (Open in ParaView is still there).

\paragraph{The crash, and why it was not in the source.} Every \code{Fem::FemPostPipeline} created in the GUI
crashed in \code{SoFCColorGradient::setMarkerLabel}. The stack (reproduced with a three-line script, \code{lldb} on
the stopped frame) ends in \code{strlen} on text that is not a pointer. Reading the registers showed
\code{setMarkerLabel} stepping through the label array in strides of 144 bytes while the array holds 128-byte
elements. Cause: the environment moved to Coin 4.0.10 on 25 Oct, which shrank \code{sizeof(SbString)} from 144 to
128 bytes (\code{CC\_STRING\_MIN\_SIZE} became $128-\mathrm{sizeof(char*)}-\mathrm{sizeof(size\_t)}$), but conda keeps
the original modification time of a package's headers (22 June), so ninja never recompiled the 414 objects that
include Inventor headers and were older than the update. They went on using the old layout against the new library.
Nothing in \code{SoFCColorGradient} or the FEM view provider was wrong, and no Vector Labs commit touched them. The
fix was to delete those objects and let ninja recompile them (about 25\,min on 18 cores, no source change), then
\file{package/bruceworx/rebuild-stale-objects.sh <package> [-f]} lists or deletes the objects older than any conda
package so the same thing cannot hide after the next \code{pixi} update. \textbf{The staged Vector Labs.app was
built from these stale objects too}; it needs one rebuild from this tree to get the fix. The crash reports in
\file{DiagnosticReports} from the app (\code{vectorlabs-2026-10-31-*.ips}, \code{Keystone-*.ips}) show other stacks
(TechDraw \code{QGSPage::findParent}, a Python-callback in Coin's dragger, \code{DocInfo::slotDeleteDocument}); the
same stale layouts are a plausible cause of some of them, but I did not test that.

\paragraph{Decided while building.}
\begin{itemize}[nosep]
  \item \textbf{Scale the file, not the pipeline.} \code{FemPostPipeline.scale()} leaves an unstructured grid empty
        in this build (\code{writeVTK}: ``Empty data object'', even for a factor of 1), so
        \code{scaled\_copy} writes an \code{mm} copy of the .vtu/.vtp next to the original with the points scaled
        (vtk and numpy, in place) and FEM reads that. The cause of the empty grid in \code{scale()} is not found.
  \item \textbf{One pipeline per part patch} for the pressure, named for the part, so each can be hidden or coloured on
        its own; the plane is a clip filter of the volume pipeline (the headless path already made it).
  \item The workaround in \code{bruceworx\_sim.\_no\_post\_pipeline} (the static-study result import) is left as it is:
        it was a choice not to make a pipeline there, and its comment still names the crash. It can go when you want
        FEM's own post objects beside a static result.
\end{itemize}

\paragraph{Tests.} GUI, \file{TestBruceWorxCfd.py} (32 tests, 0 failures): a bare \code{Fem::FemPostPipeline} read
from a .vtu in the window does not crash; the whole results view from a synthetic case (a Delaunay mesh with $p$ and
$U$ plus two boundary patches, no OpenFOAM): one surface for the part and none for the inlet, the volume pipeline
hidden, the surface coloured by $p$, the cut by $U$ with the Magnitude component. The end-to-end OpenFOAM tests now
expect the pipeline in both the headless and the GUI run. Headless \code{-t TestBruceWorxCfd}: OK (4 skipped).

\paragraph{Open questions for you.}
\begin{enumerate}[nosep]
  \item \textbf{Rebuild the app now?} The running app has the crash for every FEM post object until it is rebuilt
        from this tree (the stale-object fix is in the build directory, not in a commit).
  \item \textbf{Streamlines}: write our own (vtk's \code{vtkStreamTracer} on the mm copy, shown as a Part-less
        Coin node or an exported .vtp pipeline), or leave them to ParaView?
  \item \textbf{Surface pressure scale}: the colour bar autoscales to the part's range. Pin it to $\pm$ the dynamic
        pressure $\tfrac12\rho U^2$ so two runs compare by eye?
\end{enumerate}
